% GENERATED FILE. Edit canonical lessons, then run npm run generate:books.
% canonical-source-hash: fnv1a64:13a1c8c8d830585f
% canonical-lessons: KA-C234-sibbandi, KA-C234-naukara, KA-C234-hajari, KA-C234-pali, KA-C234-mungada, KA-R234-first-pass-words-from-chapters-230-231, KA-R234-second-pass-words-from-chapters-232-234

\chapter{Staff, Employee, Attendance, Work shift, Advance}
\label{ch:ka-staff-employee-attendance-work-shift-advance}
\hlchaptermodality{234}

\begin{chapteropening}
\textbf{By the end of this chapter:} \emph{I can say staff, an employee, attendance, a work shift and an advance.}

Complete the last lesson of chapter 234: I can say staff, an employee, attendance, a work shift and an advance.
\end{chapteropening}

\section[\texorpdfstring{\kn{ಸಿಬ್ಬಂದಿ} (sibbandi)}{ಸಿಬ್ಬಂದಿ (sibbandi)}]{\kn{ಸಿಬ್ಬಂದಿ} (sibbandi) — staff, personnel}

\label{lesson:KA-C234-sibbandi}



\begin{quote}
Before the new one: say the Kannada for science, then the Kannada for history.
\end{quote}

\subsection*{You'll want to know: \kn{ಸಿಬ್ಬಂದಿ}}
\textbf{\kn{ಸಿಬ್ಬಂದಿ}} — \emph{sibbandi} — \textquotedblleft{}staff, personnel\textquotedblright{}.

\begin{cousinweb}
\textbf{\kn{ಸಿಬ್ಬಂದಿ}} (\emph{sibbandi}) is a loan from Persian \emph{sibandī}, \textquotedblleft{}a body of paid men\textquotedblright{}, which came through Hindustani.
\end{cousinweb}

\begin{grammarlens}[title={What you've built}]
One more word for this chapter.
\end{grammarlens}

\subsection*{Your turn}
\begin{itemize}
\raggedright
  \item \emph{Say it:} \emph{sibbandi}
  \item \emph{Say it:} \emph{sibbandi}, once more
  \item \emph{Say it:} say \emph{sibbandi}
  \item \emph{Recall:} say the Kannada for science, then the Kannada for history, then say \emph{sibbandi} again
\end{itemize}

\begin{tcolorbox}[breakable,colback=teal!4,colframe=teal!35!black,title={Before you move on}]
What does \kn{ಸಿಬ್ಬಂದಿ} mean? (\textquotedblleft{}Staff, personnel\textquotedblright{}.) Say it once more.
\end{tcolorbox}
\section[\texorpdfstring{\kn{ನೌಕರ} (naukara)}{ನೌಕರ (naukara)}]{\kn{ನೌಕರ} (naukara) — an employee}

\label{lesson:KA-C234-naukara}



\begin{quote}
Before the new one: say the Kannada for history, then the Kannada for staff.
\end{quote}

\subsection*{You'll want to know: \kn{ನೌಕರ}}
\textbf{\kn{ನೌಕರ}} — \emph{naukara} — \textquotedblleft{}an employee\textquotedblright{}.

\begin{cousinweb}
\textbf{\kn{ನೌಕರ}} (\emph{naukara}) is a loan from Persian \emph{naukar}, \textquotedblleft{}servant\textquotedblright{}, which came through Hindustani.
\end{cousinweb}

\begin{grammarlens}[title={What you've built}]
One more word for this chapter.
\end{grammarlens}

\subsection*{Your turn}
\begin{itemize}
\raggedright
  \item \emph{Say it:} \emph{naukara}
  \item \emph{Say it:} \emph{naukara}, once more
  \item \emph{Say it:} say \emph{naukara}
  \item \emph{Recall:} say the Kannada for history, then the Kannada for staff, then say \emph{naukara} again
\end{itemize}

\begin{tcolorbox}[breakable,colback=teal!4,colframe=teal!35!black,title={Before you move on}]
What does \kn{ನೌಕರ} mean? (\textquotedblleft{}An employee\textquotedblright{}.) Say it once more.
\end{tcolorbox}
\section[\texorpdfstring{\kn{ಹಾಜರಿ} (hājari)}{ಹಾಜರಿ (hājari)}]{\kn{ಹಾಜರಿ} (hājari) — attendance}

\label{lesson:KA-C234-hajari}



\begin{quote}
Before the new one: say the Kannada for staff, then the Kannada for an employee.
\end{quote}

\subsection*{You'll want to know: \kn{ಹಾಜರಿ}}
\textbf{\kn{ಹಾಜರಿ}} — \emph{hājari} — \textquotedblleft{}attendance\textquotedblright{}.

\begin{cousinweb}
\textbf{\kn{ಹಾಜರಿ}} (\emph{hājari}) came through Hindustani \emph{hāzirī}, \textquotedblleft{}attendance\textquotedblright{}, from Arabic \emph{ḥāḍir}, \textquotedblleft{}present\textquotedblright{}.
\end{cousinweb}

\begin{grammarlens}[title={What you've built}]
One more word for this chapter.
\end{grammarlens}

\subsection*{Your turn}
\begin{itemize}
\raggedright
  \item \emph{Say it:} \emph{hājari}
  \item \emph{Say it:} \emph{hājari}, once more
  \item \emph{Say it:} say \emph{hājari}
  \item \emph{Recall:} say the Kannada for staff, then the Kannada for an employee, then say \emph{hājari} again
\end{itemize}

\begin{tcolorbox}[breakable,colback=teal!4,colframe=teal!35!black,title={Before you move on}]
What does \kn{ಹಾಜರಿ} mean? (\textquotedblleft{}Attendance\textquotedblright{}.) Say it once more.
\end{tcolorbox}
\section[\texorpdfstring{\kn{ಪಾಳಿ} (pāḷi)}{ಪಾಳಿ (pāḷi)}]{\kn{ಪಾಳಿ} (pāḷi) — a work shift}

\label{lesson:KA-C234-pali}



\begin{quote}
Before the new one: say the Kannada for an employee, then the Kannada for attendance.
\end{quote}

\subsection*{You'll want to know: \kn{ಪಾಳಿ}}
\textbf{\kn{ಪಾಳಿ}} — \emph{pāḷi} — \textquotedblleft{}a work shift\textquotedblright{}.

\begin{cousinweb}
\textbf{\kn{ಪಾಳಿ}} (\emph{pāḷi}) is native Dravidian, one word with Telugu \emph{pāḍi} (DEDR 4113). DEDR gives the older Kannada \emph{pāḻi}, \textquotedblleft{}row, regular order, method\textquotedblright{}: a shift is one's turn in the order. DEDR 4113 compares Sanskrit \emph{pāli}.
\end{cousinweb}

\begin{grammarlens}[title={What you've built}]
One more word for this chapter.
\end{grammarlens}

\subsection*{Your turn}
\begin{itemize}
\raggedright
  \item \emph{Say it:} \emph{pāḷi}
  \item \emph{Say it:} \emph{pāḷi}, once more
  \item \emph{Say it:} say \emph{pāḷi}
  \item \emph{Recall:} say the Kannada for an employee, then the Kannada for attendance, then say \emph{pāḷi} again
\end{itemize}

\begin{tcolorbox}[breakable,colback=teal!4,colframe=teal!35!black,title={Before you move on}]
What does \kn{ಪಾಳಿ} mean? (\textquotedblleft{}A work shift\textquotedblright{}.) Say it once more.
\end{tcolorbox}
\section[\texorpdfstring{\kn{ಮುಂಗಡ} (muṅgaḍa)}{ಮುಂಗಡ (muṅgaḍa)}]{\kn{ಮುಂಗಡ} (muṅgaḍa) — an advance (payment)}

\label{lesson:KA-C234-mungada}



\begin{quote}
Before the new one: say the Kannada for attendance, then the Kannada for a work shift.
\end{quote}

\subsection*{You'll want to know: \kn{ಮುಂಗಡ}}
\textbf{\kn{ಮುಂಗಡ}} — \emph{muṅgaḍa} — \textquotedblleft{}an advance (payment)\textquotedblright{}.

\begin{cousinweb}
\textbf{\kn{ಮುಂಗಡ}} (\emph{muṅgaḍa}) is native Dravidian: \emph{mun-}, \textquotedblleft{}before, in front\textquotedblright{} (DEDR 5020), and \emph{kaḍa}, \textquotedblleft{}debt, loan\textquotedblright{} (DEDR 1113), with \emph{k} softened to \emph{g}: money paid before.
\end{cousinweb}

\begin{grammarlens}[title={What you've built}]
One more word for this chapter.
\end{grammarlens}

\subsection*{Your turn}
\begin{itemize}
\raggedright
  \item \emph{Say it:} \emph{muṅgaḍa}
  \item \emph{Say it:} \emph{muṅgaḍa}, once more
  \item \emph{Say it:} say \emph{muṅgaḍa}
  \item \emph{Recall:} say the Kannada for attendance, then the Kannada for a work shift, then say \emph{muṅgaḍa} again
\end{itemize}

\begin{tcolorbox}[breakable,colback=teal!4,colframe=teal!35!black,title={Before you move on}]
What does \kn{ಮುಂಗಡ} mean? (\textquotedblleft{}An advance (payment)\textquotedblright{}.) Say it once more.
\end{tcolorbox}
\section[(dialogue)]{Words from chapters 230-231 — a first pass}

\label{lesson:KA-R234-first-pass-words-from-chapters-230-231}



\begin{quote}
Say the words for a work shift and an advance.
\end{quote}

\subsection*{Your turn}
Say these aloud:

\begin{itemize}
\raggedright
  \item a line, a money box, a suitcase, a hoe, a sickle
  \item a colleague, a meeting, an interview, a clerk, an officer
\end{itemize}

\begin{tcolorbox}[breakable,colback=teal!4,colframe=teal!35!black,title={Before you move on}]
A line? (\textbf{\kn{ಗೆರೆ}}, \emph{gere}.) A money box? (\textbf{\kn{ಹುಂಡಿ}}, \emph{huṇḍi}.) A suitcase? (\textbf{\kn{ಸೂಟ್ಕೇಸ್}}, \emph{sūṭkēs}.) A hoe? (\textbf{\kn{ಗುದ್ದಲಿ}}, \emph{guddali}.) A sickle? (\textbf{\kn{ಕುಡುಗೋಲು}}, \emph{kuḍugōlu}.) A colleague? (\textbf{\kn{ಸಹೋದ್ಯೋಗಿ}}, \emph{sahōdyōgi}.) A meeting? (\textbf{\kn{ಸಭೆ}}, \emph{sabhe}.) An interview? (\textbf{\kn{ಸಂದರ್ಶನ}}, \emph{sandarśana}.) A clerk? (\textbf{\kn{ಗುಮಾಸ್ತ}}, \emph{gumāsta}.) An officer? (\textbf{\kn{ಅಧಿಕಾರಿ}}, \emph{adhikāri}.)
\end{tcolorbox}
\section[(dialogue)]{Words from chapters 232-234 — a second pass}

\label{lesson:KA-R234-second-pass-words-from-chapters-232-234}



\begin{quote}
Say the words for an expense and insurance.
\end{quote}

\subsection*{Your turn}
Say these aloud:

\begin{itemize}
\raggedright
  \item an expense, insurance, a queue, experience, a resignation
  \item groceries, profit, a loss, science, history
  \item staff, an employee, attendance, a work shift, an advance
\end{itemize}

\begin{tcolorbox}[breakable,colback=teal!4,colframe=teal!35!black,title={Before you move on}]
An expense? (\textbf{\kn{ಖರ್ಚು}}, \emph{kharcu}.) Insurance? (\textbf{\kn{ವಿಮೆ}}, \emph{vime}.) A queue? (\textbf{\kn{ಸಾಲು}}, \emph{sālu}.) Experience? (\textbf{\kn{ಅನುಭವ}}, \emph{anubhava}.) A resignation? (\textbf{\kn{ರಾಜೀನಾಮೆ}}, \emph{rājīnāme}.) Groceries? (\textbf{\kn{ದಿನಸಿ}}, \emph{dinasi}.) Profit? (\textbf{\kn{ಲಾಭ}}, \emph{lābha}.) A loss? (\textbf{\kn{ನಷ್ಟ}}, \emph{naṣṭa}.) Science? (\textbf{\kn{ವಿಜ್ಞಾನ}}, \emph{vijñāna}.) History? (\textbf{\kn{ಇತಿಹಾಸ}}, \emph{itihāsa}.) Staff? (\textbf{\kn{ಸಿಬ್ಬಂದಿ}}, \emph{sibbandi}.) An employee? (\textbf{\kn{ನೌಕರ}}, \emph{naukara}.) Attendance? (\textbf{\kn{ಹಾಜರಿ}}, \emph{hājari}.) A work shift? (\textbf{\kn{ಪಾಳಿ}}, \emph{pāḷi}.) An advance? (\textbf{\kn{ಮುಂಗಡ}}, \emph{muṅgaḍa}.)
\end{tcolorbox}
